% !TEX root = ../main.tex

\begin{acknowledgements}

四年的本科生活很快就要结束了。回头看在中国科学技术大学学习和生活的这段时间，我在课程学习、工程实践和对自己兴趣的认识上都有了很多变化。本论文能够完成，离不开老师、同学、朋友和家人的帮助。借这个机会，向所有支持过我的人表示真诚的感谢。

首先，衷心感谢我的导师邢凯教授在毕业论文选题、系统设计、实验思路和论文写作过程中给予的指导。从最初确定研究方向，到后面梳理技术路线、完善系统实现、分析实验结果，邢老师都给了我很多耐心而具体的建议。老师严谨的治学态度和清晰的问题意识，也让我在完成论文的过程中受益很多。

感谢计算机科学与技术学院各位老师在本科阶段对我的培养。课堂学习为我打下了计算机系统、算法、人工智能和软件工程等方面的基础，也让我逐渐学会把一个想法落实成可以运行的系统。这些课程训练和实践经历，是我完成本论文的重要基础。

感谢身边的同学和朋友在学习、调试、讨论和生活中给我的支持。做 MarketMind 系统和写论文的过程中，我遇到过很多细碎但很耗时间的问题，包括数据处理、模型训练、实验评估和论文排版。和大家交流的过程帮我理清了不少思路，也让我在压力比较大的时候没有一直卡在原地。

感谢我的家人一直以来的理解、支持和鼓励。你们给了我稳定而温暖的后盾，让我可以专注于学习和探索。无论事情进展顺利还是遇到困难，家人的信任都是我继续往前走的重要力量。

最后，感谢中国科学技术大学给予我的学习平台和成长环境。本科四年的经历让我更加明确自己对计算机技术、人工智能应用和真实世界问题解决的兴趣。毕业论文是本科阶段的一次总结，也会是之后继续学习和探索的起点。

\end{acknowledgements}
